\section*{Chapitre \ref{ch:iv2:etfs-futures-and-index-arbitrage} --- ETF, contrats à terme et arbitrage d'indice}
\markboth{\protect\omTocMark{}{Réponses~: ETF et arbitrage d'indice}}{\protect\omTocMark{}{Réponses~: ETF et arbitrage d'indice}}

\begin{solution}{iq:iv2:etfs-futures-and-index-arbitrage:1}
La valeur liquidative\index{valeur liquidative} est de 40.00~\$ par action, et la prime est de 25 points de base.

Les avoirs valent $50 \times 20 + 30 \times 50 + 10 \times 100 + 500 = 1\,000 + 1\,500 + 1\,000 + 500 = 4\,000$, pour 100 actions de l'ETF. Et $0.10/40 = 0.25~\%$.

\iqwrong{Oublier la trésorerie. Elle appartient aux porteurs de parts, et elle représente souvent une part importante de l'écart pour un fonds proche d'une date de dividende.}
\end{solution}

\begin{solution}{iq:iv2:etfs-futures-and-index-arbitrage:2}
Il n'y a pas d'arbitrage entre 39.96~\$ et 40.06~\$. À un cours acheteur de 40.09~\$, on vend l'ETF, on achète le panier et on procède à la création, et on gagne 3~cents par action.

Pour créer, on paie $\$40.00 + \$0.06$, donc une vente au-dessus de 40.06~\$ est profitable. Pour racheter, on achète l'ETF, puis on vend le panier pour 40.00~\$ moins 0.04~\$, donc un achat en dessous de 39.96~\$ est profitable. La fourchette de juste prix de l'ETF est donc plus large que la valeur liquidative du fonds, et elle est fixée par les coûts de création et de rachat\index{création et rachat}.

\iqwrong{Appliquer le même coût des deux côtés, ou comparer le cours acheteur au cours vendeur. On vend au cours acheteur. Et la création et le rachat ont chacun leurs propres coûts, parfois très différents.}
\end{solution}

\begin{solution}{iq:iv2:etfs-futures-and-index-arbitrage:3}
Le juste prix est $100 \times 1.02 - 1.01 = \$100.99$. À 102~\$, on vend le contrat à terme et on achète l'action avec de l'argent emprunté.

On emprunte 100~\$ pour six mois afin d'acheter l'action, qui coûte 102~\$ à l'échéance. Le dividende de 1~\$, versé dans trois mois, rapporte $1~\% $ pour le trimestre restant, soit 1.01~\$ à l'échéance. On livre l'action à 102~\$, on rembourse 102~\$ et on garde les 1.01~\$. Le profit verrouillé est donc de 1.01~\$, soit l'écart entre 102~\$ et 100.99~\$.

\iqwrong{Utiliser $S(1 + rT)$ en ignorant le dividende. Le détenteur du contrat à terme ne reçoit aucun dividende, donc le contrat à terme est moins cher de la valeur du dividende à la livraison.}
\end{solution}

\begin{solution}{iq:iv2:etfs-futures-and-index-arbitrage:4}
Le notionnel d'un contrat est de 200\,000~\$. Il faut 50 contrats pour un bêta de 1, et 60 pour un bêta de 1.2.

$4\,000 \times \$50 = \$200\,000$. Puis $\$10\,000\,000/\$200\,000 = 50$. Un bêta de 1.2 signifie que le portefeuille évolue comme 12~millions de dollars de l'indice, soit 60 contrats.

\iqwrong{Couvrir le notionnel lorsque le bêta diffère de 1. Avec 50 contrats, un bêta de 1.2 laisse une exposition de marché de 2~millions de dollars.}
\end{solution}

\begin{solution}{iq:iv2:etfs-futures-and-index-arbitrage:5}
Le taux de financement implicite est de 3~\% par an. Avec un taux d'emprunt de 2.5~\%, on achète le panier avec de l'argent emprunté et on vend le contrat à terme.

La prime du contrat à terme est de $20/4000 = 0.5\%$ pour un trimestre, soit 2~\% par an, et $r - q = 2\%$, donc $r = 3\%$. C'est le cash and carry\index{cash and carry}~: on prête son argent au taux implicite de 3~\%, et on l'emprunte à 2.5~\%, soit un gain de 0.5~\% par an avant frais.

\iqwrong{Comparer les 2~\% de la prime à 2.5~\%, et conclure que le contrat à terme est bon marché. Le détenteur du panier perçoit aussi le dividende, il faut donc l'ajouter.}
\end{solution}

\begin{solution}{iq:iv2:etfs-futures-and-index-arbitrage:6}
Le contrat à terme implique 20~points d'indice de dividendes, soit un rendement de 2~\% par an.

La fair value sans dividendes est $4000 \times (1 + 5\% \times \tfrac14) = 4050$. Le contrat à terme à 4\,030 est inférieur de 20~points, c'est donc le dividende implicite\index{dividende implicite}. Et $20/4000 = 0.5\%$ pour un trimestre.

\iqwrong{Prendre chaque écart entre le spot et les contrats à terme pour des dividendes. Cet écart intègre aussi tout spread de financement, et en période de tensions il peut en contenir beaucoup.}
\end{solution}

\begin{solution}{iq:iv2:etfs-futures-and-index-arbitrage:7}
L'indice rapporte $-0.04~\%$, le fonds rapporte $-0.36~\%$, et trois fois le rendement de l'indice vaut $-0.12~\%$.

L'indice vaut $1.02 \times 0.98 = 0.9996$. Le fonds vaut $1.06 \times 0.94 = 0.9964$. Le multiple de chaque journée est exact, mais sur deux jours le rendement se compose différemment, et il est plus faible à cause de la volatilité supplémentaire.

\iqwrong{Penser qu'un fonds à effet triple rapporte trois fois le rendement sur n'importe quel horizon. Il ne promet que le multiple quotidien.}
\end{solution}

\begin{solution}{iq:iv2:etfs-futures-and-index-arbitrage:8}
Le fonds perd environ 1.82~\%. Avec la position longue sur l'indice pour ces deux jours, on perd environ 0.91~\% sur les 2~\$ investis.

Une hausse de 10~\% signifie que le fonds perd 10~\%. Puis l'indice baisse de $10/110 \approx 9.09~\%$, et le fonds gagne 9.09~\%. Donc $0.9 \times 1.0909 \approx 0.9818$. Le rendement de la position sur l'indice est nul, donc avec 1~\$ dans chaque, on termine à 1.9818~\$.

La couverture d'un jour est exacte, mais elle cesse de l'être dès que les deux positions dérivent, car l'exposition du fonds est remise à zéro chaque jour et l'on n'a pas fait de même pour sa propre position.

\iqwrong{Traiter le fonds inverse comme une couverture de long terme. Pour qu'il reste une couverture, il faut rééquilibrer entre les deux, aussi souvent que le fonds se remet à zéro.}
\end{solution}

\begin{solution}{iq:iv2:etfs-futures-and-index-arbitrage:9}
On vend 10~millions de dollars du gagnant et on achète 10~millions de dollars du perdant, soit un taux de rotation à sens unique de 5~\%. Un fonds pondéré par la capitalisation n'a besoin d'aucun échange, car ses poids évoluent seuls avec les cours.

Les avoirs deviennent 110~millions et 90~millions de dollars. La pondération égale exige 100~millions de dollars de chacun. Le total est de 200~millions de dollars, donc 10~millions représentent 5~\%.

Ainsi, la pondération par la capitalisation a un faible taux de rotation, tandis que la pondération égale vend systématiquement les gagnants et achète les perdants, et paie le coût de transaction pour le faire.

\iqwrong{Supposer qu'un fonds pondéré par la capitalisation doit trader après un mouvement de cours. Il trade sur les changements de composition de l'indice ou du nombre d'actions.}
\end{solution}

\begin{solution}{iq:iv2:etfs-futures-and-index-arbitrage:10}
La perte est de 50\,000~\$, et ces 50\,000~\$ sont débités en espèces aujourd'hui. La marge initiale est de 100\,000~\$ à l'ancien cours et de 97\,500~\$ au nouveau, donc il faut disposer de 47\,500~\$ de plus que les 50\,000~\$ restant sur le compte.

Le notionnel de la position est $10 \times 50 \times 4000 = \$2\,000\,000$, et 5~\% de ce montant font 100\,000~\$. Une baisse de 100~points sur 10 contrats à 50~\$ représente 50\,000~\$. Donc, après le règlement-livraison, le compte dispose de 50\,000~\$, et la nouvelle exigence de $5~\% \times \$1\,950\,000 = \$97\,500$ requiert 47\,500~\$ de plus.

Un contrat à terme de gré à gré de même exposition ne ferait aucun mouvement de trésorerie avant le règlement-livraison. Son gain ou sa perte serait porté au titre du risque de contrepartie. Les contrats à terme transforment le risque de marché en risque de liquidité, et il faut disposer de la trésorerie aujourd'hui.

\iqwrong{Penser qu'une position couverte ne peut pas causer de problème. Une couverture par contrats à terme d'une position longue illiquide peut donner lieu à un appel de marge en espèces alors que la position longue ne peut pas être vendue.}
\end{solution}

\begin{solution}{iq:iv2:etfs-futures-and-index-arbitrage:11}
Le fonds perd environ 2~\% par mois à cause du roll, soit une perte composée d'environ 21~\% par an sur une courbe inchangée, avant les intérêts sur sa trésorerie. Un spot inchangé ne donne pas un résultat inchangé, car le fonds continue d'acheter un contrat à terme plus cher, qui baisse ensuite vers le spot.

Chaque mois, il vend à 70 et achète à 71.40, donc il détient $70/71.40 = 0.9804$ fois le nombre de contrats. En un an, cela fait $0.9804^{12} \approx 0.788$, soit une perte de 21.2~\%.

Cette courbe est en contango\index{contango}~: le contrat à terme est au-dessus du spot. Cette prime rémunère le stockage et le financement. Le détenteur d'un contrat ne profite pas du tout de la détention du pétrole, et le coût lui est donc facturé.

\iqwrong{Dire que le fonds perd sur la transaction à chaque roll parce qu'il «~paie 1.40~». Le roll lui-même se fait au prix du marché. La perte provient de la baisse du contrat plus cher vers le spot au cours du mois.}
\end{solution}

\begin{solution}{iq:iv2:etfs-futures-and-index-arbitrage:12}
Avant la remise à zéro, les actifs valent 1.03~milliard de dollars et l'exposition vaut 3.03~milliards de dollars. Il doit donc acheter 60~millions de dollars à la clôture. Le flux vaut $L(L - 1)rA$. Un fonds à moins trois fois d'une taille identique achète 120~millions de dollars.

Son exposition était de 3~milliards de dollars, et avec une hausse de 1~\% elle vaut 3.03~milliards de dollars. Ses actifs ont augmenté de 30~millions de dollars. La cible est $3 \times 1.03 = 3.09$ milliards de dollars, donc il achète 60~millions de dollars.

En général, l'exposition après le mouvement vaut $LA(1 + r)$, les actifs valent $A(1 + Lr)$, et la cible vaut $LA(1 + Lr)$. L'écart vaut $LAr(L - 1)$.

Pour $L = -3$, $L(L - 1) = 12$, donc il trade $12 \times 1~\% \times \$1$ milliard~: il doit aussi acheter lorsque le marché monte. Vérification pour $L = 1$~: rien à faire, et c'est bien le cas.

\iqwrong{Penser qu'un fonds vendeur du marché vend lorsque le marché monte. La hausse de l'indice a accru son exposition vendeuse et réduit ses actifs, il doit donc acheter.}
\end{solution}

\begin{solution}{iq:iv2:etfs-futures-and-index-arbitrage:13}
Sur une durée $T$, la valeur du fonds vaut $(S_T/S_0)^L\exp\left(-\tfrac12L(L - 1)\sigma^2T\right)$. Avec une volatilité de 30~\% et un indice plat, le fonds à trois fois rapporte environ $-23.7~\%$, et le fonds à moins un environ $-8.6~\%$.

Le log-rendement d'un fonds remis à zéro chaque jour vaut à peu près $L$ fois le log-rendement de l'indice, avec le frein de volatilité $-\tfrac12L^2\sigma^2T$ du fonds contre $-\tfrac12\sigma^2T$ pour l'indice. La différence vaut donc $-\tfrac12(L^2 - L)\sigma^2T$.

Pour $L = 3$~: $e^{-3 \times 0.09} = e^{-0.27} \approx 0.763$. Pour $L = -1$~: $L(L - 1) = 2$, donc $e^{-0.09} \approx 0.914$.

\iqwrong{Considérer cela comme des frais ou comme une erreur du fonds. C'est l'arithmétique du rééquilibrage, et le fonds fait exactement ce qu'il annonce.}
\iqfollow{Et si l'indice a connu une tendance~?}{Une tendance peut permettre au fonds de dépasser $L$ fois le rendement~: si l'indice monte de 50~\% avec peu de volatilité, le fonds à trois fois gagne près de $1.5^3 - 1 = 237~\%$. Cela dépend du chemin suivi par l'indice.}
\end{solution}

\begin{solution}{iq:iv2:etfs-futures-and-index-arbitrage:14}
Le nouveau diviseur vaut $130/60 \approx 2.167$. Après une hausse de 10~\%, l'indice vaut 66, et il aurait valu 66 sans la division. La division n'a donc aucun effet sur le niveau, mais elle modifie les pondérations.

La division ramène la première action à 50~\$, donc la somme des cours vaut $50 + 50 + 30 = 130$. L'indice doit rester à 60, donc le diviseur vaut $130/60$. Si chaque action monte de 10~\%, la somme vaut 143, et $143/2.167 = 66$. Sans division, la somme vaut 198, et $198/3 = 66$.

Si seule la première action monte de 10~\% après la division, l'indice gagne $5/2.167 \approx 2.31$, alors qu'avant la division il aurait gagné $10/3 \approx 3.33$. La division a donc réduit le poids de cette action.

\iqwrong{Penser que la division laisse la sensibilité de l'indice inchangée. Dans un indice pondéré par les cours, les pondérations sont les cours des actions, donc une division réduit l'importance de cette action.}
\end{solution}

\begin{solution}{iq:iv2:etfs-futures-and-index-arbitrage:15}
L'indice baisse de 0.8 point d'indice. Un indice de rendement total ne baisse pas. La fair value du contrat à terme ne bouge pas ce jour-là, car le dividende figurait déjà dans les prévisions.

$4000 \times 2~\% \times 1~\% = 0.8$. Un indice de prix ne réinvestit pas, donc il baisse du montant du dividende. Un indice de rendement total réintègre le dividende, donc il reste inchangé. La fair value du contrat à terme est le spot, diminué des dividendes à venir. Le spot baisse de 0.8, et les dividendes restant à venir baissent de 0.8 également, donc les deux s'annulent, à l'intérêt près.

\iqwrong{Considérer la baisse à la date ex comme une perte. Le détenteur reçoit le cash, donc la position totale ne change pas.}
\end{solution}

\begin{solution}{iq:iv2:etfs-futures-and-index-arbitrage:16}
On gagne 20 points. De l'autre côté se trouve celui qui détient le contrat à terme en position longue~; le plus souvent, il s'agit d'un dealer de produits structurés qui se décharge de l'exposition longue aux dividendes que porte son livre. Pour la seule exposition, on négocierait un contrat à terme sur dividendes\index{contrat à terme sur dividendes} ou un swap de dividendes.

On reçoit $2.5~\%$ de 4\,000 au lieu de 2~\%, soit 100 points au lieu de 80. Le cours du contrat à terme n'a pas changé, donc les 20 points constituent le profit.

Les dealers qui vendent des produits structurés à des investisseurs reçoivent souvent des dividendes en position longue comme sous-produit, et les revendent. Les dividendes implicites peuvent donc rester sous le niveau attendu pendant de longues périodes. Les intérêts sur le dividende ont été ignorés ici.

\iqwrong{Considérer un basis trade comme sans risque. La prévision des dividendes est l'un de ses principaux risques, et elle compte davantage à longue échéance.}
\end{solution}

\begin{solution}{iq:iv2:etfs-futures-and-index-arbitrage:17}
Sur un an, le contrat à terme coûte 19 points de base et l'ETF 8. Ils sont égaux à un mois, donc en deçà d'un mois le contrat à terme est moins cher.

Le contrat à terme coûte $19H$ pour une durée de détention de $H$ années. L'ETF coûte $7H + 1$. En égalant les deux, on obtient $12H = 1$. Sur un horizon court, c'est la fourchette fixe de l'ETF qui compte, et sur un horizon long le coût courant.

\iqwrong{Comparer seulement les frais. Le taux de financement implicite est caché dans le cours du contrat à terme, et il varie. En période de stress, il peut monter bien au-dessus de 15 points de base.}
\end{solution}

\begin{solution}{iq:iv2:etfs-futures-and-index-arbitrage:18}
La fair value du spread vaut environ 30 points, donc 4\,050 est juste. À 4\,060, le contrat le plus éloigné est cher d'environ 10 points. Puis, si l'on doit rester en position longue, on peut acheter le spot via le panier à l'échéance, ou garder le contrat proche jusqu'à son échéance et passer par une autre voie, puisque le roll coûte 10 points de plus, soit environ 0.25~\%.

Le spread vaut $4020 \times (5~\% - 2~\%) \times \tfrac14 \approx 30.15$. Un calendar spread est le prix de marché du financement pour ce trimestre.

Si l'on détient le contrat proche, on peut vendre le spread (vendre le proche, acheter le lointain) à 40 points, en payant 10 points au-dessus de la fair value. On peut aussi chercher une autre exposition~: détenir jusqu'à l'échéance, puis acheter le panier ou un ETF, avec ses propres coûts. Si ceux-ci sont inférieurs à 10 points, on évite le roll.

\iqwrong{Considérer le roll comme un événement sans coût en fin de mois. Tout le monde fait son roll pendant les mêmes quelques jours, donc le calendar spread peut porter une prime.}
\end{solution}

\begin{solution}{iq:iv2:etfs-futures-and-index-arbitrage:19}
L'écart espéré vaut $-0.16~\%$ par an. Le tracking error est la variation de 0.05~\%.

Les frais coûtent 0.20~\%, et le prêt de titres rapporte 0.08~\%. Une détention de cash de 1~\% une année où l'indice bat le cash de 4~\% coûte 0.04~\%. Donc $-0.20 + 0.08 - 0.04 = -0.16$.

Tracking difference\index{tracking difference} est l'écart moyen, c'est-à-dire ce que l'on perd réellement. Le tracking error en est la volatilité. Un fonds peut afficher un faible tracking error et un mauvais tracking difference.

\iqwrong{Prendre les frais pour le coût total de détention d'un fonds. Le prêt de titres peut couvrir plus que les frais, et l'effet de frein lié au cash s'inverse sur un marché baissier.}
\end{solution}

\begin{solution}{iq:iv2:etfs-futures-and-index-arbitrage:20}
La bande d'absence d'arbitrage s'étend de 4\,040 à 4\,064.

Au-dessus de la bande, on vend le contrat à terme et on achète le panier avec de l'argent emprunté~: $4000 \times (1 + (5~\% - 2~\%) \times \tfrac12) = 4060$, et avec des coûts de 4 points, cela donne 4\,064. En dessous de la bande, on achète le contrat à terme et on vend le panier à découvert, et l'on prête le produit de la vente au taux de 4.5~\%, avec des frais d'emprunt de 0.3~\% à verser. Donc $4000 \times (1 + (4.5~\% - 0.3~\% - 2~\%) \times \tfrac12) = 4044$, et avec les coûts, 4\,040.

Elle n'est pas symétrique, car les deux côtés font face à des taux différents et seul le côté vendeur paie des frais. Sa largeur, 24 points, se décompose en deux parties~: 8 provenant des coûts, et 16 provenant de l'écart entre les taux de financement. Ainsi, une firme disposant d'un meilleur financement a une bande plus étroite et peut faire de l'arbitrage là où d'autres ne le peuvent pas.

\iqwrong{Penser qu'un contrat à terme n'a qu'une seule fair value. Il y en a une par coût de financement, donc un contrat à terme peut se situer entre les bandes de plusieurs firmes.}
\end{solution}

\begin{solution}{iq:iv2:etfs-futures-and-index-arbitrage:21}
On gagnerait 0.30 par action de l'ETF si tout se produisait en même temps. Mais la vente prend trois jours, avec un écart type d'environ 0.87, donc il ne s'agit pas d'arbitrage. Le 95 est la meilleure estimation. Le 100 publié est périmé.

Acheter à 95 et recevoir des obligations valant 95.50, moins 0.20, laisse 0.30. Le risque vaut $0.5~\% \times \sqrt 3 \approx 0.87~\%$. Un écart type représente donc presque trois fois le profit.

Les obligations se négocient rarement. La valeur liquidative est établie à partir de cours périmés ou de modèles, alors que l'ETF se négocie en continu. En période de stress, l'ETF est souvent la plateforme de formation des prix, et la «~décote~» correspond largement à la perte réelle, pas encore reflétée dans les valorisations. La décote d'un ETF ne représente au plus qu'un arbitrage modeste.

\iqwrong{Dire que l'ETF «~a décroché~» à cause d'une décote. La décote reflète des obligations qui n'ont pas pu être vendues à leur valorisation.}
\end{solution}

\begin{solution}{iq:iv2:etfs-futures-and-index-arbitrage:22}
On gagne environ 0.8 cent net par action, soit 800~\$ sur 100\,000 actions. Un mouvement défavorable de la base de 1 point de base représente 0.4 cent par action, soit 400~\$.

On gagne 1 cent par rapport à la fair value à l'exécution. La couverture coûte la moitié de la fourchette du contrat à terme~: $0.5 \text{ pb} \times \$40 = 0.2$ cents. Donc $1 - 0.2 = 0.8$.

Le risque est l'écart entre l'ETF et le contrat à terme~: le contrat à terme peut ne pas répliquer le panier de l'ETF, et sa valeur intègre des hypothèses de financement et de dividendes. On détient donc un risque supplémentaire tant qu'on ne l'a pas compensé. La moitié du profit peut disparaître sur un mouvement de la base de 1 point de base.

\iqwrong{Considérer la couverture comme parfaite. Le profit d'un teneur de marché tient largement à la petitesse de ce risque par rapport à la demi-fourchette.}
\end{solution}

\begin{solution}{iq:iv2:etfs-futures-and-index-arbitrage:23}
Le teneur de marché n'a à couvrir dans le panier que le net, soit 10~millions de dollars, ce qui coûte environ 2\,000~\$ au total. Par dollar d'ETF échangé, sur 190~millions de dollars, cela représente environ 0.1 point de base. Ainsi, la fourchette d'un ETF peut être plus étroite que celle du panier, car la plupart des transactions sont compensées par d'autres investisseurs.

$10$ millions de dollars à 2 points de base font 2\,000~\$. Sur 190~millions de dollars, cela représente environ 0.105 point de base. Les transactions de compensation ne passent jamais par le panier, donc il ne faut qu'une faible couverture pour cette part. Le reste de la fourchette rémunère le risque.

\iqwrong{Supposer qu'un ETF ne peut pas être moins cher que ce qu'il détient. Sa liquidité dépasse la somme de ses parties tant que les échanges se font dans les deux sens. Sur un marché à sens unique, la fourchette revient vers celle du panier.}
\end{solution}

\begin{solution}{iq:iv2:etfs-futures-and-index-arbitrage:24}
Le swap est moins cher de 10 points de base par an~: le cash coûte le financement plus 60, et le swap coûte le financement plus 50.

En cash, on paie 30 et on perd $15~\% \times 2~\% = 30$ sur les dividendes. Avec le swap, on paie 40 et on perd $5~\% \times 2~\% = 10$.

Avec un swap, on assume le risque de crédit du dealer, généralement réduit par la marge quotidienne, et celui-ci peut mettre fin à l'opération. On n'a aucun droit de vote, et les conditions peuvent changer. Le dealer se couvre en détenant l'action, et il peut bénéficier d'un meilleur traitement fiscal ou d'un financement plus avantageux que celui de l'investisseur~: c'est la source de l'écart. Il faut vérifier la situation juridique, car les règles fiscales applicables à ces swaps varient d'un pays à l'autre.

\iqwrong{Comparer seulement les deux spreads de financement. La fuite liée aux dividendes a ici la même taille que le financement, et c'est elle qui tranche la question.}
\end{solution}

\begin{solution}{iq:iv2:etfs-futures-and-index-arbitrage:25}
Ils doivent acheter 1.6~milliards de dollars. Cela représente 16 jours à 100~\% du volume, et 80 jours à 20~\%. La règle donne une pression sur les prix de 8~\%. Il faut s'en méfier, car 16 jours de volume sont très loin du domaine dans lequel la règle a été mesurée.

Le flottant vaut $\$10$ milliards $\times 80~\% = \$8$ milliards, et 20~\% de ce montant font 1.6~milliards de dollars. On a alors $1.6/0.1 = 16$, et $16/0.2 = 80$. La règle donne $2~\% \times \sqrt{16} = 8~\%$.

Raisons de s'en méfier~: cette demande est annoncée, ce qui attire des vendeurs qui ne traiteraient pas autrement, ainsi que des arbitragistes qui détiennent l'action à l'avance. Le volume le jour de l'effet est plusieurs fois supérieur à la normale. On ne verra donc peut-être qu'une partie de cette pression sur les prix, et une partie de celle-ci s'inversera.

\iqwrong{Dire que les fonds doivent acheter sur 80 jours. Ils veulent le cours de clôture à la date d'effet, donc ils achèteront presque tout en un seul moment.}
\end{solution}

\begin{solution}{iq:iv2:etfs-futures-and-index-arbitrage:26}
L'espérance de gain est nulle. Le gain médian est d'environ 0.41~\$, et l'on perd de l'argent environ 28~\% du temps. Une perte d'environ 2~\$ ou plus survient une fois tous les 20 ans. La décroissance est réelle pour un chemin typique, et elle est payée par une perte rare et très importante.

La valeur de chaque fonds est $\prod_i(1 \pm 3r_i)$, et avec des rendements quotidiens indépendants et de moyenne nulle, son espérance vaut exactement 1. Les positions vendeuses sur les deux fonds ont donc une espérance de gain nulle, malgré la forte décroissance de chacun. D'après la formule, avec une dérive $\exp(-\tfrac12L(L - 1)\sigma^2T)$ et $x$ le multiple de l'indice, la valeur totale des deux fonds vaut environ $e^{-0.27}x^3 + e^{-0.54}x^{-3}$. En $x = 1$, cela donne environ 1.35, donc le gain vaut environ 0.65 si l'indice termine à plat. Le gain maximal, si la volatilité réalisée vaut exactement 30~\%, est d'environ 0.67.

On perd lorsque le log-rendement de l'indice est supérieur à environ $+28~\%$ ou inférieur à environ $-37~\%$. La valeur croît alors comme le cube du gain de l'indice, et cette perte n'a pas de limite supérieure. Une simulation des rendements quotidiens donne la médiane, le seuil de 28~\% et la queue de distribution.

\iqwrong{Traiter la décroissance comme un rendement. Il s'agit d'un résultat pour un chemin typique, et sa logique est la même que celle de la vente d'un strangle~: espérance nulle, mais médiane positive. Même ce gain tient avant les frais d'emprunt, qui sont généralement élevés sur ces fonds.}
\end{solution}

\begin{solution}{iq:iv2:etfs-futures-and-index-arbitrage:27}
Chaque groupe vend 2.4~milliards de dollars, soit 4.8~milliards au total, ce qui représente 24~\% de la liquidité des 15 dernières minutes. Cela ressemble à une couverture en gamma négatif, car le montant est proportionnel au mouvement et s'échange dans le sens de celui-ci.

D'après la formule $L(L - 1)rA$. Pour $L = 3$~: $6 \times (-4~\%) \times \$10$ milliards $= -\$2.4$ milliards. Pour $L = -3$~: $12 \times (-4~\%) \times \$5$ milliards $= -\$2.4$ milliards.

Vérifions le premier par les expositions~: elle valait 30~milliards de dollars et vaut maintenant 28.8~milliards, alors que les actifs valent $\$10 \times 0.88 = \$8.8$ milliards, donc la cible est de 26.4~milliards de dollars. Cela correspond à une vente de 2.4~milliards de dollars. Vérifions le second~: elle valait $-\$15$ milliards et vaut maintenant $-\$14.4$ milliards, les actifs valent $\$5 \times 1.12 = \$5.6$ milliards, et la cible est de $-\$16.8$ milliards. Cela correspond à une vente de 2.4~milliards de dollars.

Un acteur en gamma négatif voit son delta évoluer contre lui, et il doit trader avec le marché. Ici, les fonds sont de fait en gamma négatif~: leur position par mouvement vaut $L(L - 1)A$. Une tendance en fin de journée peut donc être amplifiée par les flux des fonds eux-mêmes.

\iqwrong{Penser que les fonds long et inverses s'annulent. Ils tradent tous deux dans le sens du mouvement, donc leurs flux s'additionnent.}
\end{solution}

\begin{solution}{iq:iv2:etfs-futures-and-index-arbitrage:28}
La probabilité de l'état haussier est $5/9$. Le prix à terme est d'environ 108.35, et le prix du contrat à terme d'environ 108.75. Ils diffèrent parce que l'action est positivement corrélée aux taux d'intérêt. Si le taux n'est pas aléatoire, ou n'est pas corrélé à l'action, l'écart est nul.

L'action doit croître en moyenne au taux sans risque~: $100 \times 1.04 = 120p + 84(1 - p)$, ce qui donne $p = 20/36 = 5/9$.

L'obligation zéro-coupon à deux périodes vaut $\tfrac{1}{1.04}\left(\tfrac59/1.06 + \tfrac49/1.02\right) \approx 0.92292$, donc le prix à terme vaut $100/0.92292 \approx 108.35$. Le prix du contrat à terme est le prix final espéré sous cette probabilité risque-neutre~: $\tfrac59 \times 120 \times 1.06 + \tfrac49 \times 84 \times 1.02 \approx 108.75$.

L'acheteur d'un contrat à terme reçoit des gains quotidiens dans l'état haussier, lorsque le taux est élevé et qu'il peut les réinvestir à ce taux, et subit des pertes dans l'état baissier, lorsqu'il peut les financer à bon marché. Le détenteur a donc un avantage sur un contrat à terme ferme, et le prix du contrat à terme doit être plus élevé pour compenser cet avantage. C'est l'ajustement de convexité\index{ajustement de convexité}.

\iqwrong{Traiter les contrats à terme ferme et les contrats à terme standardisés comme identiques à tout moment. Ils sont proches pour les contrats courts sur actions, mais ils diffèrent sensiblement pour les contrats longs et pour les taux eux-mêmes.}
\end{solution}

\begin{solution}{iq:iv2:etfs-futures-and-index-arbitrage:29}
Les pondérations sont de 0.54 sur $A$ et 0.24 sur $B$, par dollar d'ETF. La volatilité résiduelle est d'environ 8.8~\%, soit une réduction de variance de 76~\%.

Les covariances entre les actions de couverture sont $0.04$, $0.03$ et $0.09$. Les covariances de l'ETF avec elles valent $0.8 \times 0.18 \times 0.2 = 0.0288$ et $0.7 \times 0.18 \times 0.3 = 0.0378$. En résolvant $\Sigma w = c$, on obtient $w = (0.54, 0.24)$. La variance expliquée vaut $w^\top c = 0.024624$ sur $0.0324$. Le reste vaut $18~\% \times \sqrt{0.24} \approx 8.8~\%$.

Davantage de titres donnent une variance plus faible dans l'échantillon, mais vos coûts et l'erreur d'estimation augmentent tous deux. Une grande matrice de covariance s'ajuste au bruit~: hors échantillon, les pondérations peuvent être instables, et la variance réelle est plus élevée que ne le laisse croire l'ajustement. Il faut donc choisir un petit nombre de couvertures stables.

\iqwrong{Utiliser les pondérations par capitalisation boursière des deux actions dans l'ETF. Une couverture de variance minimale n'est pas la même chose qu'une tranche des positions détenues.}
\end{solution}

\begin{solution}{iq:iv2:etfs-futures-and-index-arbitrage:30}
Si l'on reste vendeur, on gagne 6. Avec un retournement libre, on peut espérer 7, donc l'option vaut 1 point. Avec 0.5 point par retournement, elle vaut 0.875 point. On ne l'utilise que le jour 2, si l'écart a atteint $-2$.

Remontons en arrière. Le jour 3, l'écart passe à 0. Détenir la position le jour 2 vaut donc $|b|$~: vendeur si l'écart est positif, acheteur s'il est négatif. Le jour 1, le mouvement suivant est de moyenne nulle, donc la position n'a pas d'influence en espérance, et la valeur vaut $E|b_2|$. Depuis 10~: $(14 + 6)/2 = 10$. Depuis 2~: $(6 + 2)/2 = 4$. Depuis aujourd'hui~: $(10 + 4)/2 = 7$.

Le gain ne provient que du chemin 6, 2, $-2$, de probabilité $\tfrac14$. Rester vendeur y perd 2, tandis que passer acheteur rapporte 2. Cela fait 4 de plus, et $\tfrac14 \times 4 = 1$. Avec le coût de 0.5, cela donne $\tfrac14 \times 3.5 = 0.875$.

C'est une part réelle de la valeur d'un basis trade~: on détient une option de sortie anticipée et de retournement, et c'est l'une des raisons pour lesquelles on garde une position lorsque la base est serrée.

\iqwrong{Dire qu'une base appelée à converger ne vaut que son écart d'ouverture. Si l'on détient le droit de retourner la position, la convergence rapporte davantage que cet écart.}
\end{solution}

\begin{solution}{iq:iv2:etfs-futures-and-index-arbitrage:31}
Vous obtenez un rendement d'environ $-12.5~\%$. L'arbitrage ne résorbe pas la prime, car personne ne peut créer de nouvelles parts. Pour vendre à découvert, il faut trouver un emprunt, donc on le paie et l'on risque un short squeeze. Il faudrait vendre de petites quantités, dimensionnées pour un nouvel élargissement, et seulement avec un catalyseur qui permette à nouveau la création.

La valeur du titre vaut $1.05/1.20 = 0.875$. Le sous-jacent a monté de 5~\%, et la position perd tout de même 12.5~\%.

Normalement, une prime attire les participants autorisés, qui créent des parts en vendant l'ETF contre le panier. Si la création est suspendue, le fonds se comporte comme un fonds fermé, donc son cours peut s'écarter de sa valeur. La prime peut s'élargir avant de se résorber, et l'emprunt peut devenir cher, si bien que l'on peut avoir raison et perdre tout de même sa position.

\iqwrong{Acheter le fonds pour le mouvement du sous-jacent sans regarder la prime. On achète à la fois le sous-jacent et la prime, et celle-ci ne va pas durer.}
\end{solution}

\begin{solution}{iq:iv2:etfs-futures-and-index-arbitrage:32}
Un encours vendeur à découvert supérieur à 100~\% est possible, car une même action peut être prêtée, vendue, puis de nouveau prêtée et vendue. Les positions acheteuses totalisent 220 millions~: les 100 millions de parts en circulation plus 120 millions issus des ventes à découvert. Les teneurs de marché peuvent créer davantage de parts pour les livrer ou pour répondre à la demande d'emprunt. Si l'ETF clôture, le risque principal est que les actifs du fonds soient réalisés, et que les vendeurs à découvert doivent racheter. Les détenteurs sont payés par l'intermédiaire de leurs courtiers, les vendeurs devant couvrir ce qu'ils doivent.

Un prêteur détient une créance sur un courtier, et l'acheteur de l'autre côté compte aussi comme détenteur de l'action. Les positions acheteuses sont donc égales aux parts en circulation plus l'encours vendeur à découvert. Chaque position vendeuse est compensée par une position acheteuse supplémentaire.

Le processus de création donne aux ETF un moyen de pallier la pénurie~: un participant autorisé peut créer des parts et les livrer à ceux qui en ont besoin. Un encours vendeur élevé sur un ETF est donc souvent un outil de couverture.

Le risque ne porte généralement pas sur les actifs, qui sont détenus pour le compte du fonds. Il tient plutôt à ce qu'un agent ou une chaîne de prêteurs ne livre pas, ou qu'en période de stress, le débouclage de la chaîne prenne du temps. Les investisseurs peuvent vérifier la part des titres prêtés.

\iqwrong{Dire qu'un encours vendeur supérieur à 100~\% signifie que le fonds a des «~parts fantômes~» et qu'il est sur le point de s'effondrer. La valeur des actifs est garantie, et les positions acheteuses supplémentaires sont compensées par des positions vendeuses.}
\end{solution}

\begin{solution}{iq:iv2:etfs-futures-and-index-arbitrage:33}
L'économie attendue est de 0.2 point de base du fonds, et le risque de tracking error est d'environ 0.45 point de base du fonds. On peut raisonnablement faire l'un ou l'autre, et un fonds jugé sur son tracking error mensuel pourrait ne pas le faire.

L'action représente 0.1~\% du fonds, donc les 2~\% d'impact économisés valent 0.2 point de base. La volatilité sur cinq jours vaut $2~\% \times \sqrt 5 \approx 4.47~\%$, et pour une pondération de 0.1~\%, cela fait environ 0.45 point de base. Une position construite régulièrement sur les cinq jours exposerait à environ 0.30.

C'est un ratio de Sharpe honnête pour un titre isolé. Le problème est l'échelle~: le gain n'est réel que tant que peu d'autres font la même chose, et le mandat d'un fonds indiciel conduit souvent à renoncer à ce gain. De nombreux ajustements de ce type sur une année diversifient le risque, ce qui plaide pour les faire.

\iqwrong{Comparer 3~\% et 1~\% sans les ramener à la taille du fonds. Ces deux chiffres, rapportés au rendement du fonds, sont minimes, et ce qui compte est la fréquence à laquelle cela se produit sur l'ensemble des changements.}
\end{solution}

\begin{solution}{iq:iv2:etfs-futures-and-index-arbitrage:34}
La valeur du fonds vérifie
\[
\frac{L_T}{L_0} = \left(\frac{S_T}{S_0}\right)^{\!L}\exp\left(-(L - 1)rT - \tfrac12L(L - 1)\sigma^2T - fT\right).
\]
Pour $L = 3$, $\sigma = 20~\%$ et un indice à plat avec des taux nuls, elle vaut environ 0.887. Avec une hausse de 10~\% et des taux de 4~\%, elle vaut environ 1.090. Une option d'achat à la monnaie d'un an sur le fonds, avec des taux nuls, vaut environ 23.6 pour 100. Trois options d'achat de ce type sur l'indice valent environ 23.9.

Le fonds détient une exposition $L$ et emprunte $L - 1$ au taux $r$~: $dL_t/L_t = L\,dS_t/S_t - (L - 1)r\,dt - f\,dt$. D'après le lemme d'Ito\index{lemme d'Ito}, $d\ln S = dS/S - \tfrac12\sigma^2dt$ et $d\ln L_t = L\,dS/S - (L - 1)r\,dt - f\,dt - \tfrac12L^2\sigma^2dt$. En multipliant la première par $L$ et en soustrayant, le terme en $dS/S$ disparaît. Le fonds est donc une fonction déterministe de l'indice actuel, et il n'a aucune autre dépendance au chemin, pour une volatilité $\sigma$ constante.

Application numérique~: $e^{-3 \times 0.04} = e^{-0.12} \approx 0.887$ avec $r = 0$. Avec la hausse et le taux, $1.331 \times e^{-0.08 - 0.12} \approx 1.090$.

L'option d'achat~: le fonds a une volatilité $L\sigma = 60~\%$, donc son call vaut $100(2\Phi(0.3) - 1) \approx 23.6$. Une option d'achat sur l'indice vaut 7.97, et trois de ces options font 23.9.

Les deux valeurs sont proches, mais les profils de paiement diffèrent. Trois calls sur l'indice rapportent de manière linéaire en fonction de l'indice, alors que le call sur le fonds rapporte en fonction de son cube, à un facteur d'échelle près. Ainsi, le call sur le fonds rapporte moins si l'indice monte un peu, et il rapporte davantage pour une forte hausse.

\iqwrong{Supposer que la formule reste valable si la volatilité change. L'exposant utilise alors la variance réalisée, et le fonds dépend de la trajectoire de la volatilité. Ainsi, la valeur du fonds, et avec elle le paiement d'une option sur ce fonds, dépendent de la variance réalisée, ce qui n'est pas le cas du paiement d'une option sur l'indice.}
\end{solution}

\begin{solution}{iq:iv2:etfs-futures-and-index-arbitrage:35}
On attend environ 2~millions de dollars sur les 100~millions de dollars engagés, soit 2~\%, avec un écart type d'environ 4.5~\%, soit environ 4.5~millions de dollars. Si d'autres détiennent 400~millions de dollars, la demande nette des fonds sur le marché vaut 1.2~milliards de dollars, la pression est d'environ 6.9~\%, et la moitié de celle-ci subsiste, de sorte que le gain tombe à environ 1.5~\% avant la concurrence supplémentaire à l'entrée. On s'abstiendrait lorsque beaucoup détiennent la même position et que l'action a déjà dépassé le prix prévu par la règle.

Le pic vaut $2~\% \times \sqrt{16} = 8~\%$. La moitié était intégrée à l'annonce, donc il reste 4~\% à venir. L'entrée coûte $2~\% \times \sqrt{1} = 2~\%$~: 100~millions de dollars représentent un jour de volume. Le gain vaut donc $4~\% - 2~\% = 2~\%$, soit 2~millions de dollars.

Le risque est le mouvement de l'action sur les cinq jours, $2~\% \times \sqrt 5 \approx 4.47~\%$, soit 4.47~millions de dollars si l'on n'est pas couvert. Une couverture par le secteur élimine alors une grande part de la composante de marché. Cela représente environ 0.45 écart type par événement~: c'est une opération modeste, même si elle fonctionne.

Avec 400~millions de dollars détenus par d'autres, la demande nette vaut 1.2~milliards de dollars, et $2~\% \times \sqrt{12} \approx 6.93~\%$. Si la moitié est intégrée au prix, il reste environ 3.46~\%, à comparer au coût de 2~\% de l'entrée, qui augmente maintenant parce qu'ils achètent en même temps. Cela fait environ 1.46~\%.

On s'abstiendrait lorsque le mouvement de l'action depuis l'annonce est nettement supérieur à ce que cela indique, car les arbitragistes ont déjà encombré la position, ou parce que la demande des fonds peut être satisfaite par un gros détenteur susceptible de vendre à la clôture. Le risque de retournement est d'autant plus grand que d'autres sortent tous en même temps.

\iqwrong{Acheter à l'annonce, puisque «~les fonds doivent acheter~». Ils le doivent, et tout le monde le sait, donc l'opération peut être anticipée. Et après la clôture, le cours retombe souvent.}
\end{solution}
